\documentclass[11pt]{article}
\usepackage[T1]{fontenc}
\usepackage{lmodern}
\usepackage[margin=2.5cm]{geometry}
\usepackage{booktabs}
\usepackage{amsmath}
\usepackage[hidelinks]{hyperref}
\usepackage{xcolor}

\newcommand{\code}[1]{\texttt{#1}}
\emergencystretch=3em
\newcommand{\wk}{WUKONG}

\title{E1: Structure of graphs built by \wk, LightRAG and GraphRAG\\\large UDA-Bench NBA, experiment report}
\author{}
\date{October 9, 2026}

\begin{document}
\maketitle

\section{Goal}

We build a graph from the same corpus with three systems, using the same language model, and
measure how the resulting graphs can be navigated and queried. The statistics are chosen so that
they can be read without knowing the domain: size, node and relation vocabularies, edge
direction and properties. Recall against the benchmark's gold tables is reported only as a
control, to check that the three graphs contain comparable information.

\section{Setup}

\paragraph{Corpus and gold.} The NBA dataset of UDA-Bench: 216 English Wikipedia articles,
141 about players, 30 about teams, 29 about cities and 16 about team owners. The gold tables
(player, team, city, manager) have one row per article; the row ID is the document number.

\paragraph{Systems.} All three use \code{gpt-5.6-luna} with low reasoning effort and
\code{text-embedding-3-small}.
\begin{itemize}
  \item \textbf{\wk}, engine at commit \code{99d591e}, with the extraction schema in \code{workspace-v1/}: four entity types
    (Player, Team, City, Owner) with typed fields, and three relationship types
    (PlaysFor, with a \code{tenure} field; LocatedIn; Owns, with a \code{since\_year} field).
    Chunks of 800 tokens with 120 overlap; export to Neo4j CSV.
  \item \textbf{LightRAG} 1.5.7, default extraction, chunks of 1200 tokens with 100 overlap,
    16 concurrent LLM calls (\code{lightrag\_build.py}).
  \item \textbf{GraphRAG} 3.2.0, standard indexing with its default entity types
    (organization, person, geo, event), one gleaning, chunks of 1200 tokens with 100 overlap,
    default prompts (\code{graphrag-config/}).
\end{itemize}

\paragraph{Schema hygiene.} The first version of the \wk{} schema (\code{workspace/}) used
example values taken from the corpus, including gold answers (two players' birth dates, two
cities' states, a team's owner, and subject names). We replaced every example by a value that
is neither a gold cell nor the name of a subject entity, leaving descriptions, instructions,
formats and options unchanged (\code{workspace-v1/}, built by \code{make\_schema\_variants.py},
which rejects any schema that fails this check). The reported \wk{} graph is built with the
clean schema. Two conventions of the gold tables are used on purpose as target-schema knowledge:
the closed vocabulary \{Frontcourt, Backcourt\} for a player's position, and the date format
\code{YYYY/M/D}. The baselines receive no schema.

\paragraph{Build cost.} Table~\ref{tab:cost}. Output tokens include reasoning tokens. Runs were
executed concurrently with other runs, so wall times are indicative.

\begin{table}[h]
\centering
\begin{tabular}{lrrrr}
\toprule
 & LLM calls & Input tokens & Output tokens & Wall time \\
\midrule
\wk{} & 4,463 & 6.5M & 1.1M & 18 min \\
LightRAG & 6,341 & 17.3M & 7.8M & 92 min \\
GraphRAG & 24,965 & 35.3M & 15.1M & 101 min \\
\bottomrule
\end{tabular}
\caption{Cost of building each graph over the 216 documents. GraphRAG's calls are
3,362 graph extractions, 16,511 description summaries and 5,092 community reports; it also used
8.2M embedding tokens.}
\label{tab:cost}
\end{table}

\section{Measures}

\paragraph{Graph.} For \wk{} we count domain nodes and edges; the provenance nodes and edges that
link every object to its source chunk and document are left out. For LightRAG we read its
GraphML graph; for GraphRAG its entity and relationship tables.

\paragraph{Relation vocabulary, own labels.} The label of an edge is what the system stores:
the relationship type in \wk, the set of keywords in LightRAG. GraphRAG stores no relation
label, only a free-text description and a weight. For a labeling we report the number of
distinct labels, the fraction used by a single edge, the number of most frequent labels needed
to cover 50\% and 90\% of the edges, and the mean number of distinct labels per pair of node
types (over pairs with at least 100 edges).

\paragraph{Relation vocabulary, derived predicates.} To measure LightRAG and GraphRAG in the same
way, and to give GraphRAG a label at all, an LLM (same model and effort as the builds) reads the
description of every edge, together with the names of its two endpoints, and names the main
relation as a short predicate in lower snake case, 1--3 words, present tense, without dates,
numbers or names, and states which endpoint is its subject (\code{derive\_labels.py}; the prompt's
examples come from an unrelated domain). Edges are labeled independently, 100 per call: 980 calls,
5.5M input and 1.5M output tokens for 97,931 edges, all of which received a label.

\paragraph{Synonym merging.} To discount surface variation, all labels are embedded
(\code{text-\allowbreak embedding-\allowbreak 3-small}, 512 dimensions) and merged by leader clustering: in decreasing
order of frequency, each label joins the most similar existing leader if the cosine similarity
is at least $t$, and otherwise becomes a leader (\code{relation\_vocab.py}). We report
$t = 0.85$. At $t = 0.80$ the largest cluster already merges opposite relations
(\code{plays\_for} with \code{plays\_against}), so 0.85 is the most aggressive threshold that still
merges only synonyms; lower thresholds are in \code{relation\_vocab.json}.

\paragraph{Direction.} For a derived predicate, an edge is stored against its relation when the
subject of the predicate is the edge's target. We also report the fraction of predicates with at
least 10 edges where at least 10\% of the edges go each way. In \wk{} direction is fixed by the
endpoint types of each relationship type. LightRAG stores an undirected graph.

\paragraph{Properties.} Typed property values per node, and the fraction of nodes whose content is
only free text (a description).

\paragraph{Control: recall against gold.} Each gold row is matched to graph nodes by the
article's title and the gold name, after normalization (case, accents, punctuation), and by
aliases: for people, the same surname with a compatible first name (prefix or common nickname);
for teams, the nickname alone; for cities, \emph{City, State} and \emph{City of X}. \wk{} nodes
must also have the right type; LightRAG and GraphRAG nodes may have any type. A gold attribute
value (non-empty, and non-zero for counts) is found in \wk{} if the typed field holds it (numbers
within 1\%, city areas in either square miles or km$^2$, since the gold table mixes them). For
LightRAG and GraphRAG it is found if it appears anywhere in the descriptions of the matched nodes
or of their edges: dates in the usual written forms, counts only within 60 characters of a word
naming what is counted, positions through \emph{forward/center} and \emph{guard}. This is an upper
bound: the value is in the graph but not in a form a query can select. A gold relation (player's
team, team's city, team's owners, owner's team) is found in \wk{} if an edge of the right type
joins the matched nodes, and in LightRAG and GraphRAG if any edge joins them
(\code{e1\_analysis.py}).

\section{Results}

\begin{table}[h]
\centering
\small
\begin{tabular}{llrrr}
\toprule
 & & \wk & LightRAG & GraphRAG \\
\midrule
Size & Nodes & 2,113 & 34,162 & 37,917 \\
 & Edges & 3,376 & 44,040 & 53,891 \\
\midrule
Node vocabulary & Node types & 4 & 14 & 99 \\
 & Node types used once & 0 & 1 & 94 \\
\midrule
Relation vocabulary, & Labels & 3 & 25,151 & none \\
own labels & Labels used once & 0\% & 79\% & -- \\
 & Labels to cover 90\% of edges & 1 & 20,747 & -- \\
\midrule
Relation vocabulary, & Labels & 3 & 4,388 & 4,428 \\
derived predicates & Labels used once & 0\% & 56\% & 52\% \\
 & Labels to cover 50\% / 90\% of edges & 1 / 1 & 51 / 1,079 & 47 / 936 \\
 & Labels per node-type pair (mean) & 1 & 310 & 736 \\
\midrule
Synonyms merged & Labels & 3 & 3,488 & 3,477 \\
($\cos \ge 0.85$) & Labels to cover 90\% of edges & 1 & 738 & 633 \\
\midrule
Direction & Edges stored against their relation & 0\% & 48\%$^{a}$ & 16\% \\
 & Frequent labels used in both directions & 0\% & 85\% & 56\% \\
\midrule
Properties & Typed property values per node & 1.5 & 0 & 0 \\
 & Nodes described only by free text & 0\% & 100\% & 100\% \\
\midrule
Control: & Entities found & 214/216 & 214/216 & 214/216 \\
recall vs.\ gold & Attribute values found$^{b}$ & 859/955 & 714/955 & 729/955 \\
 & Relations found$^{c}$ & 172/220 & 195/220 & 192/220 \\
\bottomrule
\end{tabular}
\caption{Structure of the three graphs. For \wk{} the derived and merged rows repeat its own
labels. $^{a}$LightRAG's graph is undirected. $^{b}$For LightRAG and GraphRAG, a value counts if
it appears in their text (upper bound). $^{c}$For LightRAG and GraphRAG, any edge between the two
entities counts, whatever relation it describes; most of their extra relations are players' teams
outside the NBA, which the \wk{} schema does not target.}
\label{tab:e1}
\end{table}

\paragraph{Notes.}
\begin{itemize}
  \item The three graphs find the same entities and comparable numbers of facts. \wk{} finds more
    attribute values as typed fields than the baselines hold anywhere in their text; the baselines
    connect more player--team pairs, mostly with non-NBA clubs.
  \item In the derived labelings, about 50 predicates cover half of the edges and about 1,000 are
    needed for 90\%; more than half of all predicates label a single edge. Merging synonyms at
    $t = 0.85$ reduces the 90\% figure to 738 and 633.
  \item Example: in LightRAG, the 892 edges between a player and a team in the gold tables carry 676
    distinct keyword labels; \emph{played for} labels 17 of them.
\end{itemize}

\paragraph{Limitations.} Each system was run once. As an indication of run-to-run variation, three
runs of each of two \wk{} schema variants close to \code{workspace-v1} differ by at most 3 facts
in any row of the gold comparison, 9 attribute values and 3 relations in total, 25 nodes and
76 edges. Derived predicates come from labeling
each edge independently, without a shared vocabulary, which is the situation of a user who wants
to query these graphs by relation. The gold tables contain errors (mixed units for city areas,
legal names for owners), which affect all systems alike.

\section*{Reproduction}

From \code{paper/uda-nba/}: \code{run\_full.sh} builds the baseline graphs,
\code{run\_variants.sh} builds \wk{} with the clean schemas, then \code{e1\_analysis.py},
\code{derive\_labels.py} (once per baseline), \code{relation\_vocab.py} and \code{e1\_table.py},
which writes \code{runs/e1/summary.csv}. Data, environments and runs are not versioned.

\end{document}
